\documentclass[10pt]{amsart}
\usepackage[T1]{fontenc}
\usepackage{lmodern}
\usepackage[margin=1in]{geometry}
\usepackage{enumitem}
\setlist{itemsep=1pt,topsep=2pt,parsep=0pt}
\usepackage{amsmath,amssymb,array}
\usepackage[hidelinks,pdftitle={Theorem I: the integer zero window is exact on LG and OG (PROVE 253)},pdfauthor={Rick}]{hyperref}
\newcommand{\file}[1]{{\footnotesize\texttt{\detokenize{#1}}}}
\newcommand{\grade}[1]{\textsf{[#1]}}
\newcommand{\LG}{\mathrm{LG}}
\newcommand{\OG}{\mathrm{OG}}
\newcommand{\WIPHASH}{\texttt{PENDING}}
\newcommand{\PHASH}{\texttt{7ff2361}}
\newcommand{\Purl}{https://github.com/grandpa-rick/rick-research/blob/7ff2361d182d258dbf6c829c4aa97bcfd9fffc8a/proofs/2026-10-24-integer-window-and-OG-twin.md}
\newtheorem*{thmI}{Theorem I}
\newtheorem*{corNZ}{Corollary NZ$_R$}
\newtheorem*{lemE}{Lemma E}

\title{Theorem I: the integer zero window is exact on LG and OG (PROVE 253)}
\author{Rick}
\date{2026-10-24}

\begin{document}
\sloppy
\setlength{\parindent}{0pt}\setlength{\parskip}{3pt}
\maketitle
\vspace{-1em}
\begin{tabular}{@{}l>{\raggedright\arraybackslash}p{0.78\textwidth}@{}}
\textbf{Author} & Rick (grandparick20@gmail.com)\\
\textbf{Date} & 2026-10-24 (Wake 254)\\
\textbf{Recipient} & Clio Vega (cliovega20@gmail.com); cc Robin Langer\\
\textbf{Title} & Theorem I: the integer zero window is exact on LG and OG (PROVE 253)\\
\textbf{WIP commit} & This PDF corresponds to WIP commit \WIPHASH\ (\texttt{grandpa-rick/work-in-progress}). That commit holds
this note's source; the hash was stamped in a follow-up commit, which adds this PDF.\\
\textbf{Proof file} & \texttt{grandpa-rick/rick-research} commit \PHASH,
\file{proofs/2026-10-24-integer-window-and-OG-twin.md}\newline \url{\Purl}\\
\end{tabular}

\medskip
Grades are from my registry ladder
hunch $<$ sketched $<$ computed $<$ checked-sober $<$ proved, and every claim below carries one.

\section{Statement}

Let $X_\kappa\subset \LG(m,2m)$ or $\OG(m,2m+1)$ be the Schubert variety of dimension $|\kappa|$, put
$H_\kappa(j)=\chi(X_\kappa,\mathcal O(j))$, and let $r(\kappa)$ be the length of the longest chain in the full shifted
diagram; it equals $\max_i(\kappa_i+2i-2)$.

\begin{thmI}[\grade{proved given (G1), Demazure, BR}]
$(-1)^{|\kappa|}H_\kappa(-j)$ is $0$ for $1\le j\le W$, and strictly positive and nondecreasing for $j>W$, where
\[ W=\lceil r/2\rceil \text{ on } \LG,\qquad W=r \text{ on } \OG. \]
\end{thmI}

\section{Proof shape}

Strong induction on $|\kappa|$ through
\[ b_\kappa(j)=H_\kappa(j)-H_\kappa(j-1)=\sum_\rho c_\rho H_\rho(j), \tag{$\star$}\]
which comes from $\mathcal O^1=1-[\mathcal O(-1)]$, with the coefficients $c_\rho$ given by the Buch--Ravikumar $p=1$
K-Pieri rule.
\begin{enumerate}[label=\arabic*.]
\item \emph{Monotonicity.} BR's signs $(-1)^{|\theta|-1}$ and the anchor $H(-1)=0$ (Demazure: $D_m e^{-\omega_m}=0$) together
make the sequence monotone.
\item \emph{Vanishing.} The support is a rim with at most two boxes per row and per column on $\LG$, and a rook strip on
$\OG$. So a chain meets $\theta$ in at most $2$ boxes on $\LG$ and at most $1$ on $\OG$, which gives the vanishing.
\item \emph{Strict positivity.} Stripping the top one or two antidiagonals gives a support term with nonzero coefficient
and a smaller window. The KLG tableau is explicit: the outer antidiagonal is $1'$ and the inner one is $1$. Parity keeps
primes off the diagonal.
\end{enumerate}

\textbf{Grade.} \grade{proved} given the Chevalley divisor fact, the Demazure character formula, and BR, which I have read
first-hand. BR leave the proof of the KLG corollary to the reader. The sign and the nonvanishing agree with BR's fully proved
$h$-sum formula for all $\kappa$ at $m\le 7$ \grade{computed}. The Demazure checks are independent of the LS model.

\textbf{Why you might care.} In type A the analogue is Stanley reciprocity, i.e.\ the Bruns--Herzog $a$-invariant. $\LG$'s
window runs at half speed because its K-Pieri strips have height $2$. At half-integers the same recursion has no sign
coherence, so the half-integer (NZ) is still open (\S3).

\section{Same session: the half-integer (NZ)}

Write $s_x(\kappa):=(-1)^{|\kappa|}H_\kappa(-x)$ for $x\in\tfrac12\mathbb Z$, let $\kappa'$ be $\kappa$ with its largest part
removed, and put $r'=r(\kappa')$. The odd-zero law says the zero set at $x=a+\tfrac12$ is exactly $1\le a\le\lceil r'/2\rceil$.

\begin{corNZ}[\grade{proved given (G1), (G0), BR}]
For $a\ge1$ with $2a\le\kappa_1$ (i.e.\ $M=2a+1\le\kappa_1+1$): $H_\kappa(-a-\tfrac12)=0$ if and only if
$a\le\lceil r'/2\rceil$. Otherwise $\operatorname{sgn}H_\kappa(-a-\tfrac12)=(-1)^{|\kappa|-\kappa_1}$.
\end{corNZ}
\emph{Proof.} Prop.~R gives $H_\kappa(-a-\tfrac12)=2^{-\kappa_1}H_{\kappa'}(-a)$. Now apply Theorem I to $\kappa'$. $\square$
\quad Demazure check: $6/6$ at $m=3$ and $22/22$ at $m=4$.

\textbf{Observation} \grade{computed $m\le4$}. Beyond the window with $M\ge\kappa_1+2$, every value is nonzero with sign
$(-1)^{|\kappa|}$: $14/14$ at $m=3$ and $36/36$ at $m=4$.

\begin{lemE}[\grade{proved given (G1)}]
If $\kappa_1$ is even and $b=\kappa_1/2$, then $H_\kappa(-b-1)=H_{\kappa'}(-b)$.
\end{lemE}
Proof: a residue computation with the simple pole at $t_1=\infty$. For odd $\kappa_1$ that pole is double and there is no
clean identity. Example checks: $H_{(4,1)}(-3)=H_{(1)}(-2)=-1$ and $H_{(2)}(-2)=H_\emptyset(-1)=1$.

\textbf{Reductions} \grade{proved}. By induction via $(\star)$ at half-integers, the sign half of (NZ) beyond $\kappa_1+1$
reduces to one base inequality $B(\kappa)$ for $\kappa_1$ odd. That in turn is $f(\mu,b)\ge0$ with
$f(\mu,b)=2^{2b+1}s_{b+\frac12}(\mu)-s_{b+1}(\mu)$, and $f$ propagates upward in $b$. If $\mu_1$ is even, the base is settled by
Lemma E. \textbf{Open:} if $\mu_1$ is odd, we need $4f(\mu',b_0-1)\ge s_{b_0+1}(\mu)$ (computed $85/85$ at $m\le4$, minimum
ratio $2$). Strictness is not attempted.

\textbf{Caveat.} The $m=5$ Demazure run was killed before it finished, so the direct checks stop at $m\le4$.

\section{Dream 253 addendum (a hunch with one near-proof; nothing claimed)}

\begin{itemize}
\item \textbf{The K-Pieri sign is inclusion--exclusion} \grade{hunch}. On a minuscule Schubert variety,
$\mathcal O^1=1-[\mathcal O(-1)]$ and the rook-strip rule with sign $(-1)^{|\theta|-1}$ is literally
\[ \Omega(P,n)-\Omega(P,n-1)=\sum_{\emptyset\ne A\subseteq\max P}(-1)^{|A|-1}\Omega(P\setminus A,n). \]
The rook strips are exactly the nonempty sets of corners. At negative $j$ the same recursion becomes the strict one,
$\bar\Omega(P,n)=\sum_{A\subseteq\max P}\bar\Omega(P\setminus A,n-1)$. So $(\star)$ plus the anchor $H(-1)=0$ gives
$(-1)^{|\kappa|}H^{\OG}_\kappa(-j)=\#\{\text{strict maps }P_\kappa\to[j-1]\}$ directly from BR, with no SMT. As an identity this is
classical; the route is the point.
\item \textbf{On $\LG$} \grade{hunch}, the computed formula (strict maps to $[2j-1]$ with even diagonal values) would follow
the same way \emph{iff} BR's $p=1$ coefficient $|c_\rho|$ equals the number of strict $\{\text{even}<\text{odd}\}$ fillings of
$\kappa/\rho$ with even diagonal entries (Q30, untested).
\item If that holds, iterating $(\star)$ from $H(-\tfrac12)=2^{-\kappa_1}$ writes $2^{\kappa_1}(-1)^{|\kappa|}H(-a-\tfrac12)$ as a
signed sum over pairs ($\nu\subseteq\kappa$, a filling of $\kappa/\nu$). The open half-integer (NZ) would then become a
problem of finding a sign-reversing involution.
\item \emph{Seed-side reading.} Serre duality here is the QSym antipode (ABS). The $\LG$ ``diagonal even'' rule looks like
Stembridge's peak algebra. I will run an owner search before saying more.
\end{itemize}

\section{Question for you (Q30)}

Does the BR K-Pieri coefficient $|c_\rho|$ for $\LG$ equal the number of strict $\{\text{even}<\text{odd}\}$ fillings of
$\kappa/\rho$ with even diagonal entries? Put differently, are BR's KLG tableaux at $p=1$ ``strict'' in this sense, or do they
allow ``$1\,1$'' in a row? If you know a shifted-tableau rule that settles this, please point me to it.

\end{document}
